% Paper 3 — Conclusion
% Building Local-Language Economic Narrative Indices

\section{Conclusion}\label{sec:conclusion}

We presented the BENI v1 method pipeline, a reproducible workflow for transforming Bangla-language newspaper articles into economic narrative measurements. The paper contributes a harmonised database design, a weak-label and LLM-assisted annotation strategy, a frozen validation protocol, and a reproducible article-level prediction and monthly index pipeline for local-language news.

On the frozen BENI v1 validation set, the repaired TF-IDF model achieves 88.17\% accuracy and 0.823 F1 on the Economic class. That is the production baseline for the current release. The majority baseline remains much weaker on the minority class, which confirms that the index still depends on a genuinely trained classifier rather than a trivial threshold rule.

The monthly BENI v1 index now exists in frozen form with article-weighted and source-balanced variants. The prototype Potrika calibration index still reveals strong long-run co-movement with macroeconomic indicators, but those prototype results are separated from the frozen release tables. Improved classification accuracy does not automatically translate to better short-run prediction, so the next validation papers should treat the frozen index as a measurement object before using it for forecasting.

For BENI v1, the current paper now supplies the frozen empirical base for later macroeconomic validation and nowcasting work.

\subsection{Limitations}

This study has five main limitations. First, the silver standard uses LLM-generated labels rather than human expert annotations. While the self-consistency check ($\kappa = 0.48$, 96\% agreement) provides a reliability ceiling, true gold-standard validation requires human annotators. Second, the clean frozen validation set is only 186 articles (69 Economic), limiting the precision of downstream comparisons. Third, the prototype macro-validation results use the older Potrika calibration index, not the final frozen BENI v1 index. Fourth, the time series used in the calibration analysis covers only 79 months (2014--2020), which is sufficient for long-run correlation but insufficient for robust out-of-sample forecasting evaluation. Fifth, the analysis is limited to Bangladesh---generalizability to other low-resource languages and economies remains to be tested.

\subsection{Future Work}

Three directions follow directly. First, human annotation of a larger set (1,000+ articles) would enable true gold-standard validation and more precise active learning estimates. Second, validating the frozen BENI v1 index on the extended sample will test source-balanced narrative indicators against the broader macroeconomic series. Third, incorporating the frozen BENI v1 index into a nowcasting framework for Bangladesh inflation---as outlined in the broader research program---would test whether the long-run narrative signal can be combined with high-frequency official statistics for real-time economic monitoring.
